% ═══════════════════════════════════════════════════════════════════════
\section{Matrices stochastiques et états d'équilibre}
% ═══════════════════════════════════════════════════════════════════════

\begin{frame}{Définition 4.5 et proposition 4.3 --- Matrices stochastiques}
  \small Dans les deux modèles, la loi converge vers un unique équilibre,
  quel que soit le départ. Est-ce toujours le cas ?
  \begin{Enonce}{Définition 4.5 --- Matrice stochastique}
    \small Une matrice carrée $M$ est \defi{stochastique} si ses
    coefficients sont positifs ou nuls et si chaque colonne a une somme
    égale à $1$ : chaque colonne est un vecteur de probabilités, et
    $\vec e^{\mathsf T}M=\vec e^{\mathsf T}$.
  \end{Enonce}
  \small Toute matrice de transition est stochastique, et réciproquement.
  \begin{Enonce}{Proposition 4.3 --- Une matrice stochastique conserve les lois}
    \small Si $M$ est stochastique et $\vec v$ un vecteur de probabilités,
    alors $M\vec v$ (et donc $M^n\vec v$) est un vecteur de probabilités.
  \end{Enonce}
  \small Lecture matricielle :
  $\vec e^{\mathsf T}(M\vec v)=(\vec e^{\mathsf T}M)\vec v
  =\vec e^{\mathsf T}\vec v=1$ (le reste : un exercice du chapitre).
\end{frame}

\begin{frame}{Théorèmes 4.4 et 4.5 --- La valeur propre $1$}
  \begin{Enonce}{Théorème 4.4 --- Valeur propre d'une matrice stochastique}
    \small Si $M$ est stochastique, alors $1$ est une valeur propre de $M$.
  \end{Enonce}
  \begin{AuTableau}
    \footnotesize $M$ et $M^{\mathsf T}$ ont le même polynôme
    caractéristique ; et $M^{\mathsf T}\vec e=\vec e$ (transposer
    $\vec e^{\mathsf T}M=\vec e^{\mathsf T}$).
  \end{AuTableau}
  \small Ce vecteur propre a-t-il des composantes positives ?
  Perron--Frobenius (admis) :
  \begin{Enonce}{Théorème 4.5 --- Perron--Frobenius, cas stochastique}
    \small Soit $M$ stochastique. (i) Toute valeur propre $\lambda$
    (réelle ou complexe) vérifie $|\lambda|\leq1$. (ii) Il existe un
    vecteur de probabilités $\vec q$ tel que $M\vec q=\vec q$.
  \end{Enonce}
  \footnotesize Toute chaîne à un nombre fini d'états a au moins un
  équilibre.
\end{frame}

\begin{frame}{Exemples 4.4 et 4.5 --- Deux chaînes qui ne convergent pas}
  \small
  \begin{itemize}
    \item $I$ (taille $2$) : chaque état est conservé. Tout vecteur de
      probabilités est un équilibre, et $\vec x_n=\vec x_0$ : équilibre
      non unique, limite qui dépend du départ. Cela arrive dès que deux
      groupes d'états ne communiquent pas.
    \item $S=\mat{0}{1}{1}{0}$ : l'état change à chaque étape. Unique
      équilibre $\big(\tfrac12,\tfrac12\big)$, mais depuis l'état $1$ la
      loi alterne entre $(1,0)$ et $(0,1)$.
      \begin{itemize}
        \item Valeurs propres $1$ et $-1$ : la composante selon $-1$ change
          de signe au lieu de disparaître. Chaîne périodique.
        \item La proportion du temps passé dans chaque état converge,
          elle, vers $\tfrac12$.
      \end{itemize}
  \end{itemize}
  \vspace{4pt}
  Dans les deux cas, $1$ n'est pas la seule valeur propre de module $1$.
\end{frame}

\begin{frame}{Définition 4.6 --- Matrice stochastique régulière}
  \begin{Enonce}{Définition 4.6 --- Matrice stochastique régulière}
    Une matrice stochastique $M$ est \defi{régulière} s'il existe un
    entier $k\geq1$ tel que tous les coefficients de $M^k$ soient
    strictement positifs.
  \end{Enonce}
  \small
  \begin{itemize}
    \item Après $k$ étapes, tous les états sont accessibles depuis tous
      les états.
    \item Météo et migration : régulières dès $k=1$.
    \item $I$ et $S$ ne le sont pas : leurs puissances valent $I$ ou $S$.
  \end{itemize}
\end{frame}

\begin{frame}{Théorème 4.6 --- Convergence vers l'équilibre}
  \begin{Enonce}{Théorème 4.6 --- Convergence vers l'équilibre}
    \small Soit $M$ une matrice stochastique régulière de taille $N$.
    \begin{enumerate}
      \item[(i)] $1$ est une valeur propre simple de $M$, et toute autre
        valeur propre vérifie $|\lambda|<1$ ;
      \item[(ii)] il existe un unique vecteur de probabilités $\vec q$ tel
        que $M\vec q=\vec q$, à composantes strictement positives ;
      \item[(iii)] pour tout vecteur de probabilités $\vec x_0$,
        $\vec x_n=M^n\vec x_0\to\vec q$.
    \end{enumerate}
  \end{Enonce}
  \begin{AuTableau}
    \footnotesize (i), (ii) : Perron--Frobenius, admis. (iii) si $M$ est
    diagonalisable : $\vec x_0=\alpha_1\vec q+\sum\alpha_i\vec v_i$, donc
    $\vec x_n=\alpha_1\vec q+\sum\alpha_i\lambda_i^n\vec v_i\to\alpha_1\vec q$ ;
    la somme des composantes reste $1$ (prop.~4.3), d'où $\alpha_1=1$.
  \end{AuTableau}
  \small La vitesse : le plus grand $|\lambda_i|$, $i\geq2$. Météo :
  $\num{0.4}$ ; migration : $\num{0.8}$, plus lent (tableaux 4.1 et 4.2).
\end{frame}

\begin{frame}{Remarque 4.2 --- Proportion du temps et temps de retour}
  \begin{Rem}{Remarque 4.2}
    \small
    \begin{itemize}
      \item Première lecture (théorème 4.6) : $\P{X_n=i}\to q_i$.
      \item Deuxième lecture, sur une seule trajectoire : la proportion
        des étapes passées dans l'état $i$ tend vers $q_i$, avec
        probabilité $1$ --- le théorème ergodique (admis), une loi des
        grands nombres pour des étapes dépendantes. Vrai aussi pour les
        chaînes périodiques, si l'équilibre est unique.
      \item Troisième lecture : la chaîne revient en moyenne dans
        l'état $i$ toutes les $1/q_i$ étapes. Météo : un jour de pluie
        tous les $3$ jours, un beau jour tous les $\num{1.5}$ jour.
    \end{itemize}
  \end{Rem}
\end{frame}

\begin{frame}{Pour la classe 4.3 --- Simuler une chaîne de Markov en classe}
  \begin{Classe}{Simuler une chaîne de Markov en classe}
    \small
    \begin{itemize}
      \item Probabilités fréquentielles et théoriques : au cœur de
        l'enseignement dès le premier cycle.
      \item Dé à dix faces : s'il fait beau, « pluie » sur $1$ et $2$ ;
        s'il pleut, « beau » sur $1$ à $4$. Chaque équipe simule un mois ;
        on met en commun.
      \item La proportion de beaux jours se rapproche de $\tfrac23$.
        Question non évidente : pourquoi ne dépend-elle pas du premier
        jour ?
      \item Les jours de pluie arrivent en séquences : l'occasion de
        discuter la conception erronée d'un hasard « bien mélangé ».
    \end{itemize}
  \end{Classe}
\end{frame}
